% ============================================================
% Experiment 3 — Multi-Objective and Personalized Optimisation
% ============================================================

\section{Experiment 3: Multi-Objective Pareto Optimisation}
\label{sec:experiment3}

% ============================================================
\subsection{Motivation}
% ============================================================

Experiment 2 reduced the serendipity problem to a scalar: find the single
$\alpha^* = 0.35$ that maximises the harmonic mean over all users.
This is useful but hides the underlying trade-off — at $\alpha = 0.35$
some users receive high unexpectedness at the cost of lower relevance, and
others the reverse.  Multi-Objective Particle Swarm Optimisation (MOPSO)
makes this trade-off explicit by computing a Pareto front: the set of
$\alpha$ values from which no objective can be improved without worsening
the other.  Two questions are addressed here.

\begin{enumerate}
    \item \textbf{(Experiment 3a)} Does the Pareto front confirm the scalar
          optimum from Experiment 2, and what does the full trade-off curve
          look like?
    \item \textbf{(Experiment 3b)} Can serendipity be improved by assigning
          each user their own $\alpha_u$ instead of a single global value?
\end{enumerate}

% ============================================================
\subsection{Experiment 3a: Bi-Objective MOPSO}
\label{sec:exp3a}
% ============================================================

\subsubsection{Setup}

MOPSO simultaneously optimises two population-level objectives:
mean unexpectedness $\bar{U}(\alpha)$ and mean relevance $\bar{R}(\alpha)$
across 5\,000 users.  Both objectives are evaluated by interpolating the
pre-computed grid from Experiment 2 (Section~\ref{sec:experiment2}),
making each function evaluation essentially free.  MOPSO parameters follow
the default configuration (Table~\ref{tab:default_setup}) with an archive
of 40 non-dominated solutions maintained by crowding-distance selection.

\subsubsection{Results}

The MOPSO converged to a Pareto front of 40 non-dominated solutions spanning
$\alpha \in [0.00, 1.00]$, reported in Table~\ref{tab:exp3a_front} (selected
rows).

\begin{table}[h]
\centering
\caption{Selected solutions from the 2-objective Pareto front
         (unexpectedness vs.\ relevance).  Serendipity (harmonic mean)
         is shown for reference but was not an optimisation objective.
         Knee point in bold.}
\label{tab:exp3a_front}
\begin{tabular}{cccc}
\toprule
\textbf{$\alpha$} & \textbf{Avg unexpectedness} & \textbf{Avg CF score} & \textbf{Serendipity} \\
\midrule
0.000 & 0.344 & 3.082 & 0.381 \\
0.175 & 0.537 & 3.002 & 0.464 \\
0.278 & 0.645 & 2.875 & 0.489 \\
0.349 & 0.696 & 2.780 & 0.494 \\
\textbf{0.363} & \textbf{0.712} & \textbf{2.745} & \textbf{0.494} \\
0.386 & 0.743 & 2.670 & 0.490 \\
0.455 & 0.856 & 2.341 & 0.439 \\
0.580 & 0.971 & 1.863 & 0.332 \\
1.000 & 1.000 & 1.564 & 0.241 \\
\bottomrule
\end{tabular}
\end{table}

\begin{figure}[htbp]
\centering
\includegraphics[width=0.75\textwidth]{MScTemplate/figs/pso_moo_pareto.png}
\caption{Bi-objective Pareto front (unexpectedness vs.\ relevance).
         The knee point ($\alpha = 0.363$) is marked; the scalar optimum
         from Experiment 2 ($\alpha = 0.35$) lies one step to the left on
         the front.}
\label{fig:exp3a_front}
\end{figure}

\subsubsection{Interpretation}

Two findings stand out.

\paragraph{The Pareto front confirms the scalar optimum.}
The solution at $\alpha = 0.349$ on the front achieves serendipity $= 0.494$,
identical to the Experiment 2 grid search optimum.  MOPSO, searching the full
objective space without a pre-specified weighting, recovers the same operating
point.  This cross-validates the scalar result from two independent
perspectives: a stochastic swarm (PSO in Experiment 2) and a Pareto-dominance
search (MOPSO here).

\paragraph{The knee point lies at $\alpha = 0.363$.}
The knee — the solution with the largest crowding distance, representing the
best balance between the two objectives — is at $\alpha = 0.363$, very close
to $\alpha^* = 0.35$.  Moving left from the knee (lower $\alpha$) raises
relevance only marginally while losing significant unexpectedness; moving right
gains little additional unexpectedness at a steep relevance cost.  This
confirms that $\alpha \approx 0.35$--$0.36$ is a structurally stable operating
point, not a grid artefact.

% ============================================================
\subsection{Experiment 3b: Personalized MOPSO}
\label{sec:exp3b}
% ============================================================

\subsubsection{Motivation}

Experiment 3a uses a single $\alpha$ applied identically to all users.
Users differ, however: some tolerate more unexpectedness, others prefer
relevance.  A natural extension is to assign each user their own $\alpha_u$,
turning the 1-dimensional search into a 500-dimensional one (one $\alpha_u$
per user in the sample).

\subsubsection{Setup}

The search space is $\boldsymbol{\alpha} \in [0,1]^{500}$, with the aggregate
objective vector evaluated as $(\bar{U}(\boldsymbol{\alpha}),
\bar{R}(\boldsymbol{\alpha}))$ over the 500-user sample.  Per-user serendipity
is computed directly from the pre-loaded recommendation lists, so function
evaluations remain fast.

Crucially, particles are \emph{warm-started} near the known global optimum:
\begin{equation}
    \alpha_{u,0} \sim \mathcal{N}(0.35,\; 0.15^2), \quad
    \text{clipped to } [0, 1]
    \label{eq:warmstart}
\end{equation}
A preliminary run with uniform random initialisation (v1) converged to
$\bar{\alpha} \approx 0.18$ and achieved serendipity only $0.425$ ---
well below the global optimum --- confirming that random initialisation
cannot efficiently explore a 500-dimensional space without directional
guidance.  Warm-starting near the known good region (Eq.~\ref{eq:warmstart})
addresses this.  Other parameters: 100 particles, 500 iterations, archive
of 80 solutions.

\subsubsection{Results}

Table~\ref{tab:exp3b_front} summarises the Pareto front returned by the
warm-started run.

\begin{table}[h]
\centering
\caption{Personalized MOPSO Pareto front (selected solutions).
         Mean $\alpha$ and its standard deviation across 500 users are
         shown alongside population-level objectives.}
\label{tab:exp3b_front}
\begin{tabular}{cccc}
\toprule
\textbf{Mean $\bar{\alpha}$} & \textbf{Std $\sigma_\alpha$} &
\textbf{Avg relevance} & \textbf{Avg serendipity} \\
\midrule
0.123 & 0.154 & 2.967 & 0.424 \\
0.175 & 0.151 & 2.924 & 0.440 \\
0.220 & 0.140 & 2.879 & 0.457 \\
0.281 & 0.149 & 2.779 & 0.458 \\
0.309 & 0.141 & 2.728 & 0.460 \\
0.329 & 0.131 & 2.690 & 0.461 \\
\bottomrule
\end{tabular}
\end{table}

The personalized front contains 29 non-dominated solutions.
The best serendipity achieved is $0.461$ at mean $\bar{\alpha} = 0.329$,
with $\sigma_\alpha = 0.131$ --- confirming that the algorithm assigns
meaningfully different weights across users rather than collapsing to a
single value.  The personalized front dominates 17 of the 40 solutions
on the global Pareto front from Experiment 3a.

\begin{figure}[htbp]
\centering
\includegraphics[width=0.75\textwidth]{MScTemplate/figs/pso_moo_personalized_v2.png}
\caption{Personalized MOPSO Pareto front (red) overlaid on the global
         2-objective front from Experiment 3a (blue).  The personalized
         front dominates solutions in the high-relevance region but does
         not surpass the global scalar optimum in aggregate serendipity.}
\label{fig:exp3b_front}
\end{figure}

\subsubsection{Interpretation}

\paragraph{Personalization improves the relevance–serendipity trade-off locally.}
The personalized front dominates 17/40 of the global solutions, meaning it
achieves higher relevance for the same serendipity (or higher serendipity
for the same relevance) across a substantial portion of the trade-off space.
Users who would benefit most from more unexpected recommendations can be
assigned higher $\alpha_u$, while users who need stronger relevance filtering
receive lower values.

\paragraph{Aggregate serendipity does not exceed the global optimum.}
The best aggregate serendipity on the personalized front ($0.461$) is lower
than the global scalar optimum ($0.494$ at $\alpha = 0.35$).  This is not a
failure of the approach: the bi-objective MOPSO balances serendipity and
relevance jointly across all users simultaneously, and no single operating
point is expected to dominate a carefully tuned global scalar on every metric
at once.  The result shows that personalization is most beneficial when the
goal is to accommodate diverse user preferences, not to maximise a single
population-level score.

\paragraph{The 500-dimensional search remains challenging.}
The archive filled only 29 of 80 slots, indicating partial convergence.
The warm-start (Eq.~\ref{eq:warmstart}) was essential --- the v1 run with
random initialisation failed to find the good region at all --- but even with
guidance the swarm struggles in such a high-dimensional space.  Future work
could reduce the dimensionality by clustering users and assigning one $\alpha$
per cluster rather than per user.

% ============================================================
\subsection{Conclusion}
% ============================================================

Multi-objective optimisation confirms and extends the scalar result from
Experiment 2.  The bi-objective Pareto front (Experiment 3a) identifies a
structurally stable knee at $\alpha = 0.363$, consistent with the scalar
optimum $\alpha^* = 0.35$ found by grid search and PSO.  Personalizing
$\alpha$ per user (Experiment 3b) reshapes the trade-off curve and outperforms
the global front in the high-relevance region, but does not improve aggregate
serendipity beyond the single best scalar.  Together, these results suggest
that a global $\alpha \approx 0.35$ is near-optimal for population-level
serendipity maximisation, while per-user weights offer flexibility for systems
that must serve heterogeneous user preferences.
